\section{超参数选择}

\subsection{FFN 维度比：$d_\text{ff} / d_\text{model}$}

\begin{figure}[H]
\centering
\includegraphics[width=0.95\textwidth]{slides-images/slide-033.jpg}
\caption{FFN 维度比的约定：非门控模型使用 $d_\text{ff} = 4 \times d_\text{model}$，GLU 模型使用 $d_\text{ff} \approx \frac{8}{3} d_\text{model}$\protect\footnotemark}
\end{figure}
\footnotetext{Slides 第34页。}

对于前馈层的隐藏维度 $d_\text{ff}$（相对于模型维度 $d_\text{model}$），存在一个非常稳固的约定：

\begin{itemize}
    \item \textbf{非门控 FFN}（ReLU/GELU）：$d_\text{ff} = 4 \times d_\text{model}$
    \item \textbf{GLU 变体}：$d_\text{ff} = \frac{8}{3} d_\text{model} \approx 2.67 \times d_\text{model}$
\end{itemize}

GLU 变体的缩小因子 $2/3$ 是为了保持\textbf{参数量不变}——GLU 引入了额外的门控矩阵 $V$，所以需要相应缩小 $W_1$ 和 $V$ 的维度。

\begin{figure}[H]
\centering
\includegraphics[width=0.95\textwidth]{slides-images/slide-034.jpg}
\caption{各模型的 FFN 维度比：LLaMA、Qwen、DeepSeek 等都大致遵循 $\sim$2.67 的比例\protect\footnotemark}
\end{figure}
\footnotetext{Slides 第35页。}

\begin{knowledgebox}{T5 的大胆尝试：64 倍的 FFN 比例}
Google 的 T5 11B 模型有一个非常不寻常的设置：$d_\text{model} = 1024$，$d_\text{ff} = 65536$，比例高达 \textbf{64倍}！这远超标准的 4 倍约定。T5 团队的理由是：更宽的矩阵可以获得更好的系统并行效率。但有趣的是，后续的 T5 V1.1 回到了更标准的 2.5 倍比例，说明极端偏离约定可能并不是最优选择。
\end{knowledgebox}

\begin{figure}[H]
\centering
\includegraphics[width=0.95\textwidth]{slides-images/slide-036.jpg}
\caption{Kaplan et al. 的实验：$d_\text{ff}/d_\text{model}$ 的最优比例在 1--10 之间有一个宽阔的``盆地''，4 倍是合理的选择\protect\footnotemark}
\end{figure}
\footnotetext{Slides 第37页。来自 Kaplan et al. 的 scaling laws 论文。}

\subsection{注意力头维度}

\begin{figure}[H]
\centering
\includegraphics[width=0.95\textwidth]{slides-images/slide-038.jpg}
\caption{注意力头数与维度的关系：标准做法是 $d_\text{model} = n_\text{heads} \times d_\text{head}$\protect\footnotemark}
\end{figure}
\footnotetext{Slides 第39页。}

标准约定是保持：
$$d_\text{model} = n_\text{heads} \times d_\text{head}$$

即增加头数时，\textbf{每个头的维度相应减小}，而不是增加总的注意力参数量。GPT-3、T5、PaLM、LLaMA 2 等模型的 $d_\text{model} / (n_\text{heads} \times d_\text{head})$ 比值都非常接近 1。

\begin{warningbox}{低秩瓶颈的理论担忧}
Bhojanapalli et al. (2020) 指出，当头数很多而每头维度很小时，attention 矩阵的秩会变低，可能限制模型的表达能力。但实际上，大多数采用 1:1 比例的模型表现良好，低秩瓶颈似乎在实践中不是严重问题。
\end{warningbox}

\subsection{宽深比（Aspect Ratio）}

模型的``形状''由两个关键数字决定：$d_\text{model}$（宽度）和 $n_\text{layers}$（深度）。它们的比值称为\textbf{宽深比}：

$$\text{Aspect Ratio} = \frac{d_\text{model}}{n_\text{layers}}$$

\begin{figure}[H]
\centering
\includegraphics[width=0.95\textwidth]{slides-images/slide-040.jpg}
\caption{宽深比对性能的影响：约 100（即每层约 128 个隐藏维度）是跨尺度的最优点\protect\footnotemark}
\end{figure}
\footnotetext{Slides 第41页。来自 Kaplan et al. 的 scaling laws 论文。}

\begin{importantbox}{宽深比的经验法则}
大量模型（GPT-3、LLaMA 系列）都遵循 \textbf{$d_\text{model} / n_\text{layers} \approx 128$} 的比例。Kaplan et al. 的实验表明，这个最优比例在 50M 到 1.5B 参数的不同尺度上保持一致——这是一个非常好的消息，意味着你可以在小模型上确定比例，然后直接缩放到大模型。
\end{importantbox}

\begin{knowledgebox}{宽深比与并行化策略}
宽深比不仅影响模型表达能力，还直接决定了可以使用的并行化策略：
\begin{itemize}
    \item \textbf{更深的模型}：适合流水线并行（Pipeline Parallelism），将不同层放在不同设备上
    \item \textbf{更宽的模型}：适合张量并行（Tensor Parallelism），将矩阵切分到不同设备上
\end{itemize}
张量并行需要快速网络（高带宽、低延迟），流水线并行对网络要求较低。因此，你的硬件拓扑可能反过来影响最优的宽深比选择。
\end{knowledgebox}

\begin{figure}[H]
\centering
\includegraphics[width=0.95\textwidth]{slides-images/slide-041.jpg}
\caption{Ek et al. 的研究：在训练 loss 层面，深浅不太重要（参数量决定一切）；但下游任务上，更深的模型可能更好\protect\footnotemark}
\end{figure}
\footnotetext{Slides 第42页。}

\subsection{词表大小}

\begin{figure}[H]
\centering
\includegraphics[width=0.95\textwidth]{slides-images/slide-043.jpg}
\caption{词表大小的趋势：从早期的 30K--50K 增长到现代的 100K--250K\protect\footnotemark}
\end{figure}
\footnotetext{Slides 第44页。}

词表大小有明显的增长趋势：
\begin{itemize}
    \item \textbf{早期单语模型}（GPT-2、早期 LLaMA）：30K--50K tokens
    \item \textbf{现代多语言/生产模型}（GPT-4、Qwen、Gemma）：100K--250K tokens
\end{itemize}

更大的词表能更高效地编码非英语语言和特殊字符（如 emoji），从而降低推理成本——Cohere 强调其大词表可以显著减少非英语文本的 token 数量。

\subsection{正则化：Dropout 与 Weight Decay}

\begin{figure}[H]
\centering
\includegraphics[width=0.95\textwidth]{slides-images/slide-046.jpg}
\caption{正则化方法的使用趋势：Dropout 逐渐退出，Weight Decay 持续使用\protect\footnotemark}
\end{figure}
\footnotetext{Slides 第47页。}

\begin{importantbox}{Weight Decay 的反直觉作用}
在大语言模型预训练中使用 weight decay 不是为了\textbf{防止过拟合}——模型通常只训练一个 epoch，根本不可能过拟合。Weight decay 的真正作用是\textbf{改善训练动力学}：在 cosine learning rate schedule 下降阶段，weight decay 产生一种``隐式加速''效果，使模型获得更好的训练损失。
\end{importantbox}

\begin{figure}[H]
\centering
\includegraphics[width=0.95\textwidth]{slides-images/slide-047.jpg}
\caption{Weight decay 与学习率调度的交互：在学习率下降阶段，高 weight decay 的模型快速追赶并超越\protect\footnotemark}
\end{figure}
\footnotetext{Slides 第48页。}

\begin{warningbox}{Weight decay 不影响过拟合差距}
实验表明，不同程度的 weight decay 并不改变训练 loss 和验证 loss 之间的差距。这证实了 weight decay 在预训练中的作用不是正则化，而是优化器动力学层面的——它通过与学习率调度的复杂交互，帮助模型在训练后期更快地收敛。
\end{warningbox}

\subsection{本章小结}

\begin{importantbox}{超参数选择的``安全默认值''}
\begin{enumerate}
    \item \textbf{FFN 维度比}：非门控 $4\times$，GLU $\sim$2.67$\times$
    \item \textbf{注意力头维度}：$d_\text{model} = n_\text{heads} \times d_\text{head}$（比值为 1）
    \item \textbf{宽深比}：$d_\text{model} / n_\text{layers} \approx 128$
    \item \textbf{词表大小}：100K--200K（多语言场景）
    \item \textbf{正则化}：不用 dropout，使用 weight decay
\end{enumerate}
这些默认值经过大量模型验证，可以直接采用。但也要记住：没有哪个超参数是``写在石头上''的，合理偏离也能训练出好模型。
\end{importantbox}

%% ========== Part 6: 训练稳定性 ==========

